\chapter{Universal Mealy Evaluators and Intensional Program Objects}
\label{chap:universal-mealy-evaluators}

\section{Orientation: syntax, operators, and behaviour}
\label{sec:universal-mealy-orientation}

The preceding chapters constructed the ingredients needed for a categorical account of universal
Mealy evaluation.

Chapter~3 constructed canonical residual behaviour machines.  For a fixed pair of input and output
interfaces, the behaviour assignment of a Mealy morphism is the unique strict semantic map into the
canonical behaviour object.  This is a final-semantics form of universality: it classifies residual
behaviours, not transition operators.

Chapter~5 constructed, for a carrier
\[
        I \xleftarrow{p} P \xrightarrow{q} B_0 ,
\]
the internal category
\[
        \mathrm{KlEnd}_B(P)
\]
of fibrewise output-comparison Kleisli endomorphisms of \(P\).  This internal category does not
classify behaviours.  It classifies executable transition-output operators on the fixed carrier
\(P\).

Chapter~9 constructed finite hidden execution by feedback trace.  For a sub-endospan
\[
        R\hookrightarrow m
\]
of an execution span, the associated feedback span
\[
        L_R
        =
        \begin{pmatrix}
        0 & 1\\
        1 & R
        \end{pmatrix}
\]
has trace
\[
        \Tr(L_R)\cong \mathrm{Path}(R),
\]
where
\[
        \mathrm{Path}(R)=\coprod_{n\geq 0}R^{[n]}.
\]
Thus trace supplies finite path execution.

The present chapter isolates the fixed-carrier universal-machine structure implicit in these
constructions.  The central point is that there are two distinct universal objects.

The first is extensional:
\[
        \mathsf U_B(P):=\mathrm{KlEnd}_B(P)^{\operatorname{op}}.
\]
Its arrows are not syntactic names for transition-output operators.  Its arrows are the internally
represented transition-output operators themselves, oriented as input instructions.

The second is intensional:
\[
        \mathsf{Prog}_B(P)
        :=
        \mathrm{FreeCat}\bigl(\mathsf U_B(P)_{\mathrm{gr}}\bigr).
\]
Its arrows are formal finite strings of universal instructions.  Composition is concatenation of
programs, not Kleisli composition of operators.

The canonical comparison between them is the decoding functor
\[
        \operatorname{decode}_{B,P}:
        \mathsf{Prog}_B(P)\longrightarrow \mathsf U_B(P),
\]
which sends a formal program to its compiled transition-output operator.

Thus the basic architecture of universal Mealy evaluation is
\[
\begin{tikzcd}[column sep=large,row sep=large]
\mathsf{Prog}_B(P)
\ar[r,"\operatorname{decode}_{B,P}"]
&
\mathsf U_B(P)
\ar[r,"\mathbb U_{B,P}"]
&
B .
\end{tikzcd}
\]
The first object is an intensional programming object.  The second object is an extensional
operator object.  The first map is compilation.  The second displayed Mealy morphism is evaluation.

The chapter separates three levels:
\[
\begin{array}{c|c|c}
\text{level} & \text{canonical object} & \text{equality means} \\ \hline
\text{formal programs}
    & \mathsf{Prog}_B(P),\ A,\ C
    & \text{same formal arrow in the given intensional category} \\
\text{carrier operators}
    & \mathsf U_B(P),\ \operatorname{CarOp}(\Phi)
    & \text{same represented transition-output operation on }P \\
\text{residual behaviour}
    & \Beh^{\operatorname{can}}_{A,B}
    & \text{same residual behaviour.}
\end{array}
\]

For \(\mathsf{Prog}_B(P)\), equality is equality of finite paths in the underlying graph of
\(\mathsf U_B(P)\).  Since the primitive letters of this graph are already arrows of
\(\mathsf U_B(P)\), primitive equality is already operator equality.  The intensional distinction
is therefore a distinction between formal composites and their compiled one-step operator effects.

This distinction is essential.  A Mealy input category may remember generators, labels, program
texts, derivations, circuits, proof terms, machine descriptions, or other syntactic structure.
Extensionality appears only after such syntax is interpreted as a transition-output operation on
the carrier \(P\).  The category
\[
        \mathsf U_B(P)
\]
is the canonical target of that interpretation.

The main constructions of the chapter are:
\[
\begin{array}{c|c}
\text{construction} & \text{role} \\ \hline
\mathsf U_B(P)=\mathrm{KlEnd}_B(P)^{\operatorname{op}}
    & \text{universal extensional operator category} \\
\mathbb U_{B,P}:\mathsf U_B(P)\rightsquigarrow B
    & \text{universal extensional Mealy evaluator} \\
\mathsf{Prog}_B(P)=\mathrm{FreeCat}(\mathsf U_B(P)_{\mathrm{gr}})
    & \text{canonical intensional straight-line program category} \\
\operatorname{decode}_{B,P}:\mathsf{Prog}_B(P)\to\mathsf U_B(P)
    & \text{canonical compiler} \\
\mathbb P_{B,P}:=\operatorname{decode}_{B,P}^{\ast}\mathbb U_{B,P}
    & \text{canonical intensional universal Mealy program machine.}
\end{array}
\]

The main results are:
\[
\begin{array}{c|l}
\text{genericity} &
\text{fixed-carrier Mealy structures are classified by functors into }\mathsf U_B(P);\\
\text{terminality} &
\mathbb U_{B,P}\text{ is terminal among fixed-carrier evaluators;}\\
\text{one-step full abstraction} &
\text{universal instruction equality is equality of represented carrier effect;}\\
\text{extensional collapse} &
\text{input syntax maps into }\mathsf U_B(P)\text{ through its operational image;}\\
\text{straight-line compilation} &
\mathsf{Prog}_B(P)\to\mathsf U_B(P)\text{ compiles formal programs;}\\
\text{operational quotient} &
\mathsf U_B(P)\text{ is the quotient of }\mathsf{Prog}_B(P)\text{ by operational equality;}\\
\text{diagonal obstruction} &
\text{full total extensional universality cannot generally be self-coded.}
\end{array}
\]

Classical Turing-style universality appears as an additional refinement.  It requires intensional
codes, partial evaluation, composition tracking, parameterization, and iteration.  These structures
are not the source of the universal operator category.  They describe how a chosen effective
subcategory of the universal operator category is represented by codes.

The conceptual progression is:
\[
\begin{tikzcd}[column sep=huge,row sep=large]
\text{syntax}
\ar[r,"\text{interpret}"]
&
\text{carrier operators}
\ar[r,"\text{black-box}"]
&
\text{residual behaviour}
\\
A,\ C,\ \mathsf{Prog}_B(P)
\ar[r]
&
\mathsf U_B(P)
\ar[r,dashed]
&
\Beh^{\operatorname{can}}_{A,B}.
\end{tikzcd}
\]
The dashed arrow is not a functor defined on the full universal operator category in general.  It
indicates the behavioural quotient induced by residual semantics on the operator image generated by
a given Mealy action.

\section{Ambient categorical base}
\label{sec:universal-mealy-ambient-base}

The constructions in this chapter use three kinds of structure: dependent products for
fibrewise operator representation, regular image factorizations for operational quotients, and
finite-path free internal categories for formal programs.  We isolate these requirements once.

\begin{definition}[Universal Mealy program base]
\label{def:universal-mealy-program-base}
A \textbf{universal Mealy program base} is a finitely complete regular category \(\mathcal E\)
satisfying the following conditions.

\begin{enumerate}
    \item \textbf{Fibrewise operator representation.}
    For every internal category \(B\) in \(\mathcal E\) and every carrier
    \[
            I\xleftarrow{p}P\xrightarrow{q}B_0,
    \]
    the dependent product
    \[
            \Pi_{\pi_{\mathrm{src}}}\bigl(E_B(P)^{[2]}\bigr)
    \]
    defined in Section~\ref{sec:operator-category-of-a-carrier} exists in the slice
    \(\mathcal E/(I\times I)\), and the resulting object carries the identity and Kleisli
    composition maps making \(\mathrm{KlEnd}_B(P)\) an internal category with object-of-objects
    \(I\).

    \item \textbf{Regular images of identity-on-objects internal functors.}
    If
    \[
            F:C\longrightarrow D
    \]
    is an identity-on-objects internal functor with common object-of-objects \(I\), then the
    regular image of
    \[
            F_1:C_1\longrightarrow D_1
    \]
    in the slice \(\mathcal E/(I\times I)\) is closed under identities and composition.  Hence it
    carries the unique internal category structure for which the regular epi--mono factorization
    of \(F_1\) is a factorization by identity-on-objects internal functors.

    \item \textbf{Effective operational quotients.}
    For each identity-on-objects internal functor used below, the kernel pair of its arrow map in
    \(\mathcal E/(I\times I)\) has an effective quotient, and this quotient agrees with the
    regular image factorization.

    \item \textbf{Free internal categories by finite paths.}
    For every internal graph
    \[
            G\rightrightarrows I
    \]
    used below, the finite path objects
    \[
            G_1^{[n]}
    \]
    of composable \(n\)-tuples exist, the coproduct
    \[
            \coprod_{n\geq 0}G_1^{[n]}
    \]
    exists in \(\mathcal E/(I\times I)\), and these coproducts are disjoint and stable under
    pullback.  With unit given by length-zero paths and composition given by concatenation, this
    finite-path object is the free internal category on \(G\).
\end{enumerate}
\end{definition}

Unless explicitly stated otherwise, \(\mathcal E\) denotes a universal Mealy program base.

The first clause is the only structure needed for the extensional evaluator
\[
        \mathbb U_{B,P}.
\]
The fourth clause is used only for the intensional program object
\[
        \mathsf{Prog}_B(P).
\]
The regularity and quotient clauses are used for operational images and extensional quotients.

\begin{proposition}[Locally cartesian closed bases supply operator objects]
\label{prop:lccc-operator-objects}
Let \(\mathcal E\) be locally cartesian closed.  Then, for every internal category \(B\) and every
carrier
\[
        I\xleftarrow{p}P\xrightarrow{q}B_0,
\]
the representing object
\[
        \Pi_{\pi_{\mathrm{src}}}\bigl(E_B(P)^{[2]}\bigr)
\]
exists.  Consequently the fibrewise output-comparison Kleisli endomorphisms of \(P\) are
representable.
\end{proposition}

\begin{proof}
Since \(\mathcal E\) is locally cartesian closed, pullback along every morphism has a right
adjoint between slices.  Applying this to the source-state projection
\[
        \pi_{\mathrm{src}}:P_{\mathrm{src}}\to I\times I
\]
gives the required dependent product.  The evaluation map is the counit of the adjunction
\[
        \pi_{\mathrm{src}}^{\ast}\dashv \Pi_{\pi_{\mathrm{src}}}.
\]
The identity and Kleisli composition maps are obtained by transposing, respectively, the identity
transition-output section and the sequential evaluation of two composable transition-output
sections.
\end{proof}

\section{The operator category of a carrier}
\label{sec:operator-category-of-a-carrier}

Let \(B\) be an internal category in \(\mathcal E\), and let
\[
        I \xleftarrow{p} P \xrightarrow{q} B_0
\]
be a carrier over the object \(I\) of input-state types and the object \(B_0\) of output objects.

When an input category \(A\) is fixed, we identify
\[
        I=A_0.
\]

Chapter~5 constructs the internal category
\[
        \mathrm{KlEnd}_B(P)
\]
of fibrewise output-comparison Kleisli endomorphisms of \(P\).  Its object-of-objects is \(I\).

A generalized arrow
\[
        \varphi:v\longrightarrow u
\]
of \(\mathrm{KlEnd}_B(P)\) is a family sending
\[
        x\in P_v
\]
to
\[
        \bigl(\varphi_0(x),\alpha_x^\varphi\bigr),
\]
where
\[
        \varphi_0(x)\in P_u,
        \qquad
        \alpha_x^\varphi:q(\varphi_0(x))\longrightarrow q(x)
\]
is an arrow of \(B\).

The target-to-source convention is essential.  An input arrow
\[
        a:u\longrightarrow v
\]
acts, contravariantly, as a Kleisli operation
\[
        P_v\rightsquigarrow P_u.
\]
Thus the input arrow points from \(u\) to \(v\), while the induced state operation points from the
fibre over \(v\) to the fibre over \(u\):
\[
\begin{tikzcd}[column sep=huge,row sep=large]
u \ar[r,"a"] & v \\
P_u & P_v \ar[l,swap,"\chi_P(a^{\operatorname{op}})"] .
\end{tikzcd}
\]

We recall the representing construction.  Over \(I\times I\), regard a generalized point
\[
        (v,u)
\]
as the formal arrow type
\[
        v\longrightarrow u.
\]
Let
\[
        P_{\mathrm{src}}:=\pi_1^{\ast}P,
        \qquad
        P_{\mathrm{tgt}}:=\pi_2^{\ast}P
\]
over \(I\times I\).  Thus
\[
        (P_{\mathrm{src}})_{(v,u)}=P_v,
        \qquad
        (P_{\mathrm{tgt}})_{(v,u)}=P_u.
\]
Let
\[
        \pi_{\mathrm{src}}:P_{\mathrm{src}}\longrightarrow I\times I
\]
be the source-state projection.

Define
\[
        E_B(P)^{[2]}\longrightarrow P_{\mathrm{src}}
\]
to be the object whose generalized elements over \((v,u,x)\), with \(x\in P_v\), are pairs
\[
        (y,\alpha)
\]
where
\[
        y\in P_u,
        \qquad
        \alpha:q(y)\longrightarrow q(x)
\]
is an arrow of \(B\).  Equivalently,
\[
        E_B(P)^{[2]}
        :=
        \bigl(P_{\mathrm{src}}\times_{I\times I}P_{\mathrm{tgt}}\bigr)
        \times_{B_0\times B_0}
        B_1,
\]
where \(B_1\to B_0\times B_0\) is \((s_B,t_B)\), and where
\[
        P_{\mathrm{src}}\times_{I\times I}P_{\mathrm{tgt}}
        \longrightarrow B_0\times B_0
\]
sends
\[
        (x,y)\longmapsto(q(y),q(x)).
\]
Equivalently, \(E_B(P)^{[2]}\) is defined by the pullback square
\[
\begin{tikzcd}[column sep=large,row sep=large]
E_B(P)^{[2]}
\ar[r]
\ar[d]
&
B_1
\ar[d,"{(s_B,t_B)}"]
\\
P_{\mathrm{src}}\times_{I\times I}P_{\mathrm{tgt}}
\ar[r,swap,"{(q_{\mathrm{tgt}},q_{\mathrm{src}})}"]
&
B_0\times B_0 .
\end{tikzcd}
\]
The projection
\[
        E_B(P)^{[2]}\to P_{\mathrm{src}}
\]
forgets \(y\) and \(\alpha\).

Define
\[
        \mathrm{KlEnd}_B(P)_1
        :=
        \Pi_{\pi_{\mathrm{src}}}\bigl(E_B(P)^{[2]}\bigr)
\]
as an object of the slice \(\mathcal E/(I\times I)\).  Its universal evaluation map is the counit
\[
        \operatorname{ev}:
        \mathrm{KlEnd}_B(P)_1\times_{I\times I}P_{\mathrm{src}}
        \longrightarrow
        E_B(P)^{[2]}.
\]
Thus
\[
\begin{tikzcd}[column sep=large,row sep=large]
\mathrm{KlEnd}_B(P)_1\times_{I\times I}P_{\mathrm{src}}
\ar[r,"\operatorname{ev}"]
\ar[d]
&
E_B(P)^{[2]}
\ar[d]
\\
\mathrm{KlEnd}_B(P)_1
\ar[r]
&
I\times I .
\end{tikzcd}
\]

The identities are given by
\[
        1_v(x)=(x,1_{q(x)}).
\]

Composition is Kleisli composition.  We write
\[
        \varphi\diamond\psi
\]
for the Kleisli composite that executes \(\psi\) first and \(\varphi\) second.  Thus, for
\[
        \psi:w\longrightarrow v,
        \qquad
        \varphi:v\longrightarrow u
\]
in \(\mathrm{KlEnd}_B(P)\), their composite is
\[
        \varphi\diamond\psi:w\longrightarrow u.
\]
It is obtained by first evaluating \(\psi\) and then evaluating \(\varphi\):
\[
        x\longmapsto
        \left(
            \varphi_0(\psi_0(x)),
            \alpha_x^\psi\circ\alpha_{\psi_0(x)}^\varphi
        \right).
\]
The output-comparison component is the categorical composite
\[
        q(\varphi_0(\psi_0(x)))
        \xrightarrow{\alpha_{\psi_0(x)}^\varphi}
        q(\psi_0(x))
        \xrightarrow{\alpha_x^\psi}
        q(x).
\]
In the manuscript's multiplicative convention \(b\circ a=ba\), this may equivalently be written as
\[
        \alpha_x^\psi\,\alpha_{\psi_0(x)}^\varphi .
\]

The composition pattern is:
\[
\begin{tikzcd}[column sep=huge,row sep=large]
w \ar[r,"\psi"] & v \ar[r,"\varphi"] & u
\\
P_w \ar[r,swap,"\psi"] &
P_v \ar[r,swap,"\varphi"] &
P_u .
\end{tikzcd}
\]
The lower row displays the carrier-level Kleisli execution.

The internal category
\[
        \mathrm{KlEnd}_B(P)
\]
is therefore an operator category.  Its arrows are not syntactic representatives of operators.
They are the represented fibrewise transition-output operators themselves.

\section{The extensional universal Mealy category}
\label{sec:extensional-universal-mealy-category}

\begin{definition}[Extensional universal Mealy category]
Let \(B\) be an internal category and let
\[
        I \xleftarrow{p} P \xrightarrow{q} B_0
\]
be a carrier.  The \textbf{extensional universal Mealy category} of \(P\) over \(B\) is
\[
        \mathsf U_B(P):=\mathrm{KlEnd}_B(P)^{\operatorname{op}}.
\]
\end{definition}

The opposite is forced by the target-to-source convention of Mealy execution.  An input arrow
\[
        a:u\longrightarrow v
\]
acts as an operator
\[
        P_v\rightsquigarrow P_u.
\]
Thus a Mealy structure on \(P\) is an identity-on-objects functor
\[
        A^{\operatorname{op}}\longrightarrow \mathrm{KlEnd}_B(P),
\]
or equivalently an identity-on-objects functor
\[
        A\longrightarrow
        \mathrm{KlEnd}_B(P)^{\operatorname{op}}
        =
        \mathsf U_B(P).
\]

\[
\begin{tikzcd}[column sep=huge,row sep=large]
A^{\operatorname{op}}
\ar[r,"\chi_P"]
&
\mathrm{KlEnd}_B(P)
\\
A
\ar[r,"\ulcorner -\urcorner"']
\ar[u,phantom,"\operatorname{op}" description]
&
\mathsf U_B(P)=\mathrm{KlEnd}_B(P)^{\operatorname{op}}
\ar[u,phantom,"\operatorname{op}" description] .
\end{tikzcd}
\]

\begin{definition}[Universal extensional evaluator]
The \textbf{universal extensional Mealy evaluator}
\[
        \mathbb U_{B,P}:\mathsf U_B(P)\rightsquigarrow B
\]
is the Mealy morphism whose carrier is
\[
        I\xleftarrow{p}P\xrightarrow{q}B_0
\]
and whose Kleisli action is the identity functor
\[
        1_{\mathrm{KlEnd}_B(P)}:
        \mathrm{KlEnd}_B(P)\longrightarrow \mathrm{KlEnd}_B(P).
\]
\end{definition}

In generalized elements, if
\[
        \varphi:v\longrightarrow u
        \quad\text{in}\quad
        \mathrm{KlEnd}_B(P),
\]
then the corresponding instruction is
\[
        \varphi^{\operatorname{op}}:u\longrightarrow v
        \quad\text{in}\quad
        \mathsf U_B(P).
\]
The universal evaluator executes this instruction by evaluation:
\[
        \delta_{\mathbb U}(\varphi^{\operatorname{op}},x)=\varphi_0(x),
        \qquad
        \omega_{\mathbb U}(\varphi^{\operatorname{op}},x)=\alpha_x^\varphi.
\]
The orientation is:
\[
\begin{tikzcd}[column sep=huge,row sep=large]
u
\ar[r,"{\varphi^{\operatorname{op}}\text{ in }\mathsf U_B(P)}"]
&
v
\\
P_u
&
P_v
\ar[l,swap,"{\varphi:P_v\rightsquigarrow P_u\text{ in }\mathrm{KlEnd}_B(P)}"]
\\
q(\varphi_0(x))
\ar[r,swap,"{\alpha_x^\varphi}"]
&
q(x) .
\end{tikzcd}
\]

\section{Genericity and terminality of the universal evaluator}
\label{sec:genericity-terminality-universal-evaluator}

The first universal theorem says that \(\mathbb U_{B,P}\) is the generic evaluator for all
Mealy structures on the fixed carrier \(P\).

\begin{theorem}[Genericity of the extensional universal evaluator]
\label{thm:genericity-universal-mealy}
Let \(B\) be an internal category and let
\[
        I \xleftarrow{p} P \xrightarrow{q} B_0
\]
be a carrier.  Let \(A\) be an internal category with object-of-objects \(I\).  Then the data of a
Mealy structure on the fixed carrier \(P\), with input category \(A\) and output category \(B\), is
naturally bijective to the data of an identity-on-objects internal functor
\[
        A\longrightarrow \mathsf U_B(P).
\]
Equivalently,
\[
        \mathrm{MlyStr}_{A,B}(P)
        \cong
        \mathrm{Cat}_I(A,\mathsf U_B(P)),
\]
where \(\mathrm{MlyStr}_{A,B}(P)\) denotes the collection of fixed-carrier Mealy structures on
\(P\), and \(\mathrm{Cat}_I\) denotes identity-on-objects internal functors over \(I\).

In machine form, for every Mealy morphism
\[
        \Phi:A\rightsquigarrow B
\]
with carrier \(P\), there is a unique identity-on-objects classifying functor
\[
        \ulcorner\Phi\urcorner:A\longrightarrow\mathsf U_B(P)
\]
such that \(\Phi\) is canonically identified, as a fixed-carrier Mealy structure, with
\[
        \ulcorner\Phi\urcorner^{\ast}\mathbb U_{B,P}.
\]
\end{theorem}

\begin{proof}
A Mealy structure on the fixed carrier \(P\), with input category \(A\) and output category \(B\),
is equivalently a fibrewise Kleisli action
\[
        \chi_P:A^{\operatorname{op}}\longrightarrow \mathrm{KlEnd}_B(P)
\]
which is identity on object-of-objects.

Taking opposites converts this datum into an identity-on-objects functor
\[
        \ulcorner\Phi\urcorner
        :=
        \chi_P^{\operatorname{op}}:
        A\longrightarrow
        \mathrm{KlEnd}_B(P)^{\operatorname{op}}
        =
        \mathsf U_B(P).
\]
Conversely, any identity-on-objects functor
\[
        F:A\longrightarrow \mathsf U_B(P)
\]
opposes to an identity-on-objects functor
\[
        F^{\operatorname{op}}:
        A^{\operatorname{op}}\longrightarrow \mathrm{KlEnd}_B(P),
\]
and hence determines a Mealy structure on \(P\).

The universal Mealy morphism \(\mathbb U_{B,P}\) is induced by the identity action of
\(\mathrm{KlEnd}_B(P)\).  Restricting it along
\[
        \ulcorner\Phi\urcorner:A\to\mathsf U_B(P)
\]
gives the action
\[
        A^{\operatorname{op}}
        \xrightarrow{\ulcorner\Phi\urcorner^{\operatorname{op}}}
        \mathrm{KlEnd}_B(P)
        \xrightarrow{1}
        \mathrm{KlEnd}_B(P),
\]
which is precisely \(\chi_P\).  Thus
\[
        \Phi
        =
        \ulcorner\Phi\urcorner^{\ast}\mathbb U_{B,P}
\]
as a fixed-carrier Mealy structure.  Uniqueness follows because any classifying functor with this
property must oppose to the Kleisli action \(\chi_P\).
\end{proof}

The theorem says that every fixed-carrier Mealy machine is obtained by interpreting its
instructions as arrows of
\[
        \mathsf U_B(P).
\]

The genericity statement is summarized by the pullback-of-evaluator diagram
\[
\begin{tikzcd}[column sep=large,row sep=large]
A
\ar[r,"\ulcorner\Phi\urcorner"]
\ar[dr,swap,"\Phi"]
&
\mathsf U_B(P)
\ar[d,"\mathbb U_{B,P}"]
\\
&
B .
\end{tikzcd}
\]

The same fact can be expressed as terminality.

\begin{definition}[Fixed-carrier evaluator category]
Let \(\mathrm{Eval}_B(P)\) be the category whose objects are Mealy morphisms
\[
        \Phi:A\rightsquigarrow B
\]
with carrier
\[
        I\xleftarrow{p}P\xrightarrow{q}B_0.
\]
A morphism
\[
        F:\Phi\longrightarrow\Psi
\]
from
\[
        \Phi:A\rightsquigarrow B
        \qquad\text{to}\qquad
        \Psi:C\rightsquigarrow B
\]
is an identity-on-objects internal functor
\[
        F:A\longrightarrow C
\]
such that the fixed-carrier Mealy structure of \(\Phi\) is obtained by restricting the
fixed-carrier Mealy structure of \(\Psi\) along \(F\).
\end{definition}

Thus, if
\[
        \chi_\Phi:A^{\operatorname{op}}\to\mathrm{KlEnd}_B(P),
        \qquad
        \chi_\Psi:C^{\operatorname{op}}\to\mathrm{KlEnd}_B(P)
\]
are the corresponding Kleisli actions, the condition is
\[
        \chi_\Phi
        =
        \chi_\Psi F^{\operatorname{op}}.
\]
Equivalently,
\[
\begin{tikzcd}[column sep=large,row sep=large]
A^{\operatorname{op}}
\ar[rr,"\chi_\Phi"]
\ar[dr,swap,"F^{\operatorname{op}}"]
&
&
\mathrm{KlEnd}_B(P)
\\
&
C^{\operatorname{op}}
\ar[ur,swap,"\chi_\Psi"]
&
.
\end{tikzcd}
\]

\begin{theorem}[Terminal fixed-carrier evaluator]
\label{thm:terminal-fixed-carrier-evaluator}
The universal extensional evaluator
\[
        \mathbb U_{B,P}:\mathsf U_B(P)\rightsquigarrow B
\]
is terminal in \(\mathrm{Eval}_B(P)\).
\end{theorem}

\begin{proof}
Let
\[
        \Phi:A\rightsquigarrow B
\]
be an object of \(\mathrm{Eval}_B(P)\).  By
Theorem~\ref{thm:genericity-universal-mealy}, \(\Phi\) has a unique classifying functor
\[
        \ulcorner\Phi\urcorner:A\to\mathsf U_B(P)
\]
such that
\[
        \Phi
        =
        \ulcorner\Phi\urcorner^{\ast}\mathbb U_{B,P}
\]
as a fixed-carrier Mealy structure.  This classifying functor is precisely a morphism
\[
        \Phi\longrightarrow \mathbb U_{B,P}
\]
in \(\mathrm{Eval}_B(P)\).

If
\[
        F:\Phi\to\mathbb U_{B,P}
\]
is any other morphism in \(\mathrm{Eval}_B(P)\), then the corresponding action
\[
        A^{\operatorname{op}}
        \xrightarrow{F^{\operatorname{op}}}
        \mathrm{KlEnd}_B(P)
\]
must equal \(\chi_\Phi\).  Hence
\[
        F=\ulcorner\Phi\urcorner.
\]
Therefore every object of \(\mathrm{Eval}_B(P)\) has a unique morphism to
\(\mathbb U_{B,P}\), so \(\mathbb U_{B,P}\) is terminal.
\end{proof}

\section{One-step operational full abstraction}
\label{sec:operational-full-abstraction}

The extensional universal category has an exact equality theory for represented one-step carrier
effects.  Here ``operational'' means equality of represented fibrewise transition-output sections
over the carrier \(P\), not equality after residual black-boxing.

\begin{definition}[Operational equality of universal instructions]
Let
\[
        \varphi^{\operatorname{op}},\psi^{\operatorname{op}}:u\longrightarrow v
\]
be arrows of
\[
        \mathsf U_B(P)=\mathrm{KlEnd}_B(P)^{\operatorname{op}}.
\]
They are \textbf{operationally equal on \(P\)} if their corresponding Kleisli operators
\[
        \varphi,\psi:v\longrightarrow u
\]
in \(\mathrm{KlEnd}_B(P)\) have the same represented section
\[
        P_v\longrightarrow E_B(P)^{[2]}_{(v,u)}.
\]
In internal element notation this says
\[
        \varphi_0(x)=\psi_0(x),
        \qquad
        \alpha_x^\varphi=\alpha_x^\psi
\]
for \(x\in P_v\).
\end{definition}

\begin{theorem}[One-step operational full abstraction]
\label{thm:operational-full-abstraction}
For arrows
\[
        \varphi^{\operatorname{op}},\psi^{\operatorname{op}}:u\longrightarrow v
\]
of \(\mathsf U_B(P)\), the following are equivalent:
\[
        \varphi^{\operatorname{op}}=\psi^{\operatorname{op}},
\]
and the two represented sections
\[
        P_v\longrightarrow E_B(P)^{[2]}_{(v,u)}
\]
obtained by evaluating \(\varphi\) and \(\psi\) are equal.  Equivalently,
\[
        \operatorname{ev}(\varphi,x)=\operatorname{ev}(\psi,x)
\]
for \(x\in P_v\).
\end{theorem}

\begin{proof}
Equality
\[
        \varphi^{\operatorname{op}}=\psi^{\operatorname{op}}
\]
in \(\mathsf U_B(P)\) is the same as equality
\[
        \varphi=\psi
\]
in \(\mathrm{KlEnd}_B(P)\).  The arrow object
\[
        \mathrm{KlEnd}_B(P)_1
        =
        \Pi_{\pi_{\mathrm{src}}}\bigl(E_B(P)^{[2]}\bigr)
\]
represents sections of
\[
        E_B(P)^{[2]}\to P_{\mathrm{src}}.
\]
Therefore two arrows of \(\mathrm{KlEnd}_B(P)\) are equal precisely when the corresponding
represented sections over the source-state family are equal.  These represented sections are the
evaluated maps
\[
        x\longmapsto
        \bigl(\varphi_0(x),\alpha_x^\varphi\bigr)
\]
and
\[
        x\longmapsto
        \bigl(\psi_0(x),\alpha_x^\psi\bigr).
\]
This proves the equivalence.
\end{proof}

The representation can be displayed as:
\[
\begin{tikzcd}[column sep=large,row sep=large]
P_v
\ar[r,bend left=15,"{\operatorname{ev}(\varphi,-)}"]
\ar[r,bend right=15,swap,"{\operatorname{ev}(\psi,-)}"]
&
E_B(P)^{[2]}_{(v,u)}
\ar[r]
&
B_1 .
\end{tikzcd}
\]

\begin{lemma}[Regular-epimorphic detection of operator equality]
\label{lem:regular-epi-detection-operator-equality}
For generalized families
\[
        \varphi,\psi:X\to\mathrm{KlEnd}_B(P)_1
\]
over \(I\times I\), equality \(\varphi=\psi\) is detected by equality of the corresponding
evaluated maps after a regular-epimorphic cover of the pulled-back source-state family.
\end{lemma}

\begin{proof}
Pull back
\[
        P_{\mathrm{src}}\to I\times I
\]
along
\[
        X\to I\times I
\]
to obtain the source-state family
\[
        X\times_{I\times I}P_{\mathrm{src}}\to X.
\]
The generalized families \(\varphi\) and \(\psi\) determine evaluated maps
\[
        X\times_{I\times I}P_{\mathrm{src}}
        \longrightarrow
        E_B(P)^{[2]}.
\]
If these evaluated maps agree after a regular-epimorphic cover of
\[
        X\times_{I\times I}P_{\mathrm{src}},
\]
then they agree before the cover because regular epimorphisms are epimorphisms.  By the
representing property of
\[
        \mathrm{KlEnd}_B(P)_1
        =
        \Pi_{\pi_{\mathrm{src}}}\bigl(E_B(P)^{[2]}\bigr),
\]
equality of evaluated maps is equivalent to equality of the represented generalized operators
\[
        \varphi=\psi.
\]
\end{proof}

Thus
\[
        \text{same universal instruction}
        \quad\Longleftrightarrow\quad
        \text{same represented one-step transition-output effect on }P.
\]
The category \(\mathsf U_B(P)\) is therefore fully abstract for one-step carrier-level effect.  It
does not preserve intensional distinctions between distinct syntactic expressions that denote the
same carrier operator.

\section{Extensional collapse of instruction syntax}
\label{sec:extensional-collapse-instruction-syntax}

A Mealy morphism may have an intensional input category.  Its arrows may be generators, labels,
program texts, derivations, words, proofs, or machine instructions.  The universal evaluator
interprets these arrows as operators on \(P\), and this interpretation may identify distinct
arrows.

Let
\[
        \Phi:A\rightsquigarrow B
\]
be a Mealy morphism with carrier
\[
        I\xleftarrow{p}P\xrightarrow{q}B_0.
\]
Its action is an identity-on-objects functor
\[
        \chi_\Phi:A^{\operatorname{op}}\longrightarrow\mathrm{KlEnd}_B(P).
\]
Equivalently, by Theorem~\ref{thm:genericity-universal-mealy}, it has a classifying functor
\[
        \ulcorner\Phi\urcorner:A\longrightarrow\mathsf U_B(P).
\]

\begin{definition}[Carrier operational equivalence]
Let \(a,b:u\to v\) be arrows of \(A\).  They are \textbf{carrier-operationally equivalent on \(P\)}
if their induced transition-output operators are equal:
\[
        \chi_\Phi(a^{\operatorname{op}})
        =
        \chi_\Phi(b^{\operatorname{op}})
        \quad\text{in}\quad
        \mathrm{KlEnd}_B(P).
\]
Equivalently,
\[
        \ulcorner\Phi\urcorner(a)
        =
        \ulcorner\Phi\urcorner(b)
        \quad\text{in}\quad
        \mathsf U_B(P).
\]
We write
\[
        a\equiv_P b.
\]
\end{definition}

\begin{definition}[Carrier operational image]
The \textbf{carrier operational image} of
\[
        \Phi:A\rightsquigarrow B
\]
is the identity-on-objects regular image of its action:
\[
        \operatorname{CarOp}(\Phi)
        :=
        \operatorname{RegIm}_I
        \bigl(
            \chi_\Phi:
            A^{\operatorname{op}}\to\mathrm{KlEnd}_B(P)
        \bigr).
\]
Equivalently,
\[
        \operatorname{CarOp}(\Phi)^{\operatorname{op}}
        =
        \operatorname{RegIm}_I
        \bigl(
            \ulcorner\Phi\urcorner:
            A\to\mathsf U_B(P)
        \bigr).
\]
\end{definition}

The defining factorization is
\[
\begin{tikzcd}[column sep=large,row sep=large]
A
\ar[rr,"\ulcorner\Phi\urcorner"]
\ar[dr,swap,two heads]
&
&
\mathsf U_B(P)
\\
&
\operatorname{CarOp}(\Phi)^{\operatorname{op}}
\ar[ur,swap,hook]
&
.
\end{tikzcd}
\]
On arrow objects, this is the regular epi--mono factorization in the slice over \(I\times I\).

\begin{theorem}[Extensional collapse]
\label{thm:extensional-collapse}
The classifying functor
\[
        \ulcorner\Phi\urcorner:A\longrightarrow\mathsf U_B(P)
\]
factors through the carrier-operational quotient of \(A\):
\[
        A
        \twoheadrightarrow
        \operatorname{CarOp}(\Phi)^{\operatorname{op}}
        \hookrightarrow
        \mathsf U_B(P).
\]
The kernel congruence of this factorization is precisely carrier-operational equivalence:
\[
        a\equiv_P b
        \quad\Longleftrightarrow\quad
        \ulcorner\Phi\urcorner(a)=\ulcorner\Phi\urcorner(b).
\]
\end{theorem}

\begin{proof}
By definition,
\[
        \operatorname{CarOp}(\Phi)^{\operatorname{op}}
\]
is the identity-on-objects regular image of
\[
        \ulcorner\Phi\urcorner:A\to\mathsf U_B(P).
\]
The regular image factorization of the arrow map gives
\[
        A_1
        \twoheadrightarrow
        \operatorname{CarOp}(\Phi)^{\operatorname{op}}_1
        \hookrightarrow
        \mathsf U_B(P)_1
\]
over \(I\times I\).  By the definition of a universal Mealy program base, the middle object
carries the induced internal category structure and the two displayed maps are internal functors.

The kernel congruence of the regular epimorphism identifies exactly those arrows of \(A\) with the
same image under \(\ulcorner\Phi\urcorner\).  This is carrier-operational equivalence by
definition.
\end{proof}

This theorem does not say that automata are extensional.  It says that the universal operator
evaluator receives any instruction syntax only through its operational action on the chosen carrier.
The input category \(A\) may remain intensional.  Its image in \(\mathsf U_B(P)\) is extensional.

When \(P\) is the derivative-closed behaviour carrier associated to a specification \(K\), this
construction recovers the syntactic transition-output category of Chapter~5:
\[
        \mathrm{SynTO}_{A,B}(K)
        =
        \operatorname{RegIm}_I\left(
            A^{\operatorname{op}}
            \longrightarrow
            \mathrm{KlEnd}_B(\operatorname{Der}_0(K))
        \right).
\]
Thus \(\mathrm{SynTO}_{A,B}(K)\) is the extensional operational image of the input syntax on the
minimal derivative carrier of \(K\).

\begin{proposition}[Local coding from operational density]
\label{prop:local-coding-operational-density}
Let
\[
        \mathcal C\hookrightarrow \mathrm{KlEnd}_B(P)
\]
be an identity-on-objects internal subcategory.  Suppose a Mealy morphism
\[
        \Phi:A\rightsquigarrow B
\]
with carrier \(P\) has action
\[
        \chi_\Phi:A^{\operatorname{op}}\to\mathrm{KlEnd}_B(P)
\]
which factors through \(\mathcal C\), and suppose the induced arrow map
\[
        (A^{\operatorname{op}})_1\twoheadrightarrow \mathcal C_1
\]
is a regular epimorphism in the slice over \(I\times I\).

Then every generalized family
\[
        \theta:X\to\mathcal C_1
\]
over \(I\times I\) of transition-output operators is locally named by an input arrow of \(A\).
Explicitly, there exists a regular epimorphism
\[
        e:\widetilde X\twoheadrightarrow X
\]
over \(I\times I\) and a map
\[
        \bar\theta:\widetilde X\to (A^{\operatorname{op}})_1
\]
over \(I\times I\) such that
\[
        \chi_\Phi\bar\theta=\theta e.
\]
Moreover, the evaluated transition-output operations agree after pullback along \(e\).
\end{proposition}

\begin{proof}
Pull back the regular epimorphism
\[
        (A^{\operatorname{op}})_1\twoheadrightarrow \mathcal C_1
\]
along
\[
        \theta:X\to\mathcal C_1
\]
in the slice \(\mathcal E/(I\times I)\).  This gives the pullback square
\[
\begin{tikzcd}[column sep=large,row sep=large]
\widetilde X
\ar[r,"\bar\theta"]
\ar[d,swap,"e"]
&
(A^{\operatorname{op}})_1
\ar[d,two heads]
\\
X
\ar[r,swap,"\theta"]
\ar[ur,phantom,very near start,"\lrcorner"]
&
\mathcal C_1 .
\end{tikzcd}
\]
Regular epimorphisms are pullback-stable in a regular category, so
\[
        e:\widetilde X\twoheadrightarrow X
\]
is a regular epimorphism.  The defining pullback square gives
\[
        \chi_\Phi\bar\theta=\theta e.
\]
Pulling back the universal evaluation map along these equal maps gives the same evaluated
transition-output operation over
\[
        \widetilde X\times_{I\times I}P_{\mathrm{src}}.
\]
\end{proof}

This proposition is the operational coding form of the extensional-collapse theorem.  Density of
the input syntax in an operator subcategory says that every operator in that subcategory is locally
named by syntax.

\section{The intensional universal program category}
\label{sec:intensional-universal-program-category}

The extensional universal category
\[
        \mathsf U_B(P)
\]
has all transition-output operators as primitive instructions.  Since its arrows are already
operators, its composition compiles finite instruction strings immediately.

To retain formal program structure, we forget composition and freely generate a category.

\begin{definition}[Intensional universal program category]
The \textbf{intensional universal program category} of \(P\) over \(B\) is
\[
        \mathsf{Prog}_B(P)
        :=
        \mathrm{FreeCat}\bigl(\mathsf U_B(P)_{\mathrm{gr}}\bigr),
\]
where \(\mathsf U_B(P)_{\mathrm{gr}}\) is the underlying internal graph of
\(\mathsf U_B(P)\).
\end{definition}

Thus
\[
        \mathsf{Prog}_B(P)_0=I.
\]
By the finite-path presentation of the free internal category,
\[
        \mathsf{Prog}_B(P)_1
        =
        \coprod_{n\geq 0}\mathsf U_B(P)_1^{[n]},
\]
where \(\mathsf U_B(P)_1^{[n]}\) denotes the object of composable \(n\)-tuples of arrows in the
underlying graph of \(\mathsf U_B(P)\), with
\[
        \mathsf U_B(P)_1^{[0]}=I.
\]

A generalized arrow of \(\mathsf{Prog}_B(P)\) is a finite composable word
\[
        u_0
        \xrightarrow{\varphi_1^{\operatorname{op}}}
        u_1
        \xrightarrow{\varphi_2^{\operatorname{op}}}
        \cdots
        \xrightarrow{\varphi_n^{\operatorname{op}}}
        u_n
\]
in \(\mathsf U_B(P)\).  Equivalently, each letter corresponds to a Kleisli operator
\[
        \varphi_i:u_i\longrightarrow u_{i-1}
        \quad\text{in}\quad
        \mathrm{KlEnd}_B(P).
\]

The empty word at \(u\) is the identity program at \(u\).  Composition is concatenation:
\[
[\varphi_1^{\operatorname{op}},\ldots,\varphi_m^{\operatorname{op}}]
;
[\psi_1^{\operatorname{op}},\ldots,\psi_n^{\operatorname{op}}]
=
[\varphi_1^{\operatorname{op}},\ldots,\varphi_m^{\operatorname{op}},
  \psi_1^{\operatorname{op}},\ldots,\psi_n^{\operatorname{op}}],
\]
whenever the target of the first word is the source of the second.

The distinction between free concatenation and operator composition is:
\[
\begin{tikzcd}[column sep=huge,row sep=large]
\mathsf U_B(P)_{\mathrm{gr}}
\ar[r,hook]
\ar[dr,swap,"\mathrm{id}_{\mathsf U_B(P)_{\mathrm{gr}}}"]
&
\mathsf{Prog}_B(P)
\ar[d,dashed,"\operatorname{decode}_{B,P}"]
\\
&
\mathsf U_B(P).
\end{tikzcd}
\]

The category \(\mathsf{Prog}_B(P)\) does not identify a two-step program with its compiled one-step
operator.  In general,
\[
        [\varphi^{\operatorname{op}},\psi^{\operatorname{op}}]
\]
and
\[
        [(\varphi\diamond\psi)^{\operatorname{op}}]
\]
are different arrows of \(\mathsf{Prog}_B(P)\), even though they decode to the same arrow of
\(\mathsf U_B(P)\), whenever the displayed word is composable.

\begin{proposition}[Intensionality of formal programs]
\label{prop:intensionality-formal-programs}
In the canonical finite-path construction of the free internal category, equality in
\[
        \mathsf{Prog}_B(P)
\]
is equality of formal composable paths in the underlying graph of \(\mathsf U_B(P)\).  In
particular, a one-letter word and a two-letter word are not identified merely because they decode
to the same arrow of \(\mathsf U_B(P)\).
\end{proposition}

\begin{proof}
The arrow object of the free internal category is
\[
        \mathsf{Prog}_B(P)_1
        =
        \coprod_{n\geq 0}\mathsf U_B(P)_1^{[n]}.
\]
The \(n\)-th summand records composable words of length \(n\).  Disjointness of the path-length
coproducts keeps different lengths distinct.  Composition is concatenation, not evaluation in
\(\mathsf U_B(P)\).  Hence relations coming from composition in \(\mathsf U_B(P)\) are imposed
only after applying the decoder.
\end{proof}

\begin{remark}
Since the underlying graph of \(\mathsf U_B(P)\) includes the identity arrows of
\(\mathsf U_B(P)\), the one-letter program
\[
        [1_u]
\]
is a primitive instruction of length one.  It is not the empty word at \(u\) in the free category.
The decoder identifies them:
\[
        \operatorname{decode}_{B,P}([1_u])
        =
        \operatorname{decode}_{B,P}([])
        =
        1_u.
\]
Thus even identity instructions are separated from formal identity programs before decoding.
\end{remark}

\section{Decoding and straight-line compilation}
\label{sec:decoding-straight-line-compilation}

Because \(\mathsf{Prog}_B(P)\) is free on the underlying graph of \(\mathsf U_B(P)\), the identity
graph map
\[
        \mathsf U_B(P)_{\mathrm{gr}}
        \longrightarrow
        \mathsf U_B(P)_{\mathrm{gr}}
\]
extends uniquely to an internal functor
\[
        \operatorname{decode}_{B,P}:
        \mathsf{Prog}_B(P)\longrightarrow\mathsf U_B(P).
\]

\begin{definition}[Canonical decoder]
The \textbf{canonical decoder}, or \textbf{straight-line compiler},
\[
        \operatorname{decode}_{B,P}:
        \mathsf{Prog}_B(P)\longrightarrow\mathsf U_B(P)
\]
is the unique identity-on-objects functor extending the identity map on primitive universal
instructions.
\end{definition}

In generalized elements, if
\[
        w=
        [\varphi_1^{\operatorname{op}},\ldots,\varphi_n^{\operatorname{op}}]
\]
is a word
\[
        u_0
        \xrightarrow{\varphi_1^{\operatorname{op}}}
        u_1
        \xrightarrow{\varphi_2^{\operatorname{op}}}
        \cdots
        \xrightarrow{\varphi_n^{\operatorname{op}}}
        u_n ,
\]
where
\[
        \varphi_i:u_i\to u_{i-1}
        \quad\text{in}\quad
        \mathrm{KlEnd}_B(P),
\]
then
\[
        \operatorname{decode}_{B,P}(w)
        =
        (\varphi_1\diamond\varphi_2\diamond\cdots\diamond\varphi_n)^{\operatorname{op}}.
\]
The order is important.  The composite
\[
        \varphi_1\diamond\varphi_2\diamond\cdots\diamond\varphi_n
\]
is the Kleisli operator
\[
        u_n\longrightarrow u_0
\]
which applies \(\varphi_n\) to the source state first, then \(\varphi_{n-1}\), and so on, ending
with \(\varphi_1\).

The decoder sits in
\[
\begin{tikzcd}[column sep=large,row sep=large]
\mathsf U_B(P)_{\mathrm{gr}}
\ar[r,hook]
\ar[dr,swap,"\mathrm{id}"]
&
\mathsf{Prog}_B(P)
\ar[d,"\operatorname{decode}_{B,P}"]
\\
&
\mathsf U_B(P).
\end{tikzcd}
\]

The execution order is:
\[
\begin{tikzcd}[column sep=large,row sep=large]
u_0
\ar[r,"\varphi_1^{\operatorname{op}}"]
&
u_1
\ar[r,"\varphi_2^{\operatorname{op}}"]
&
\cdots
\ar[r,"\varphi_{n-1}^{\operatorname{op}}"]
&
u_{n-1}
\ar[r,"\varphi_n^{\operatorname{op}}"]
&
u_n
\\
P_{u_0}
&
P_{u_1}
\ar[l,swap,"\varphi_1"]
&
\cdots
\ar[l,swap,"\varphi_2"]
&
P_{u_{n-1}}
\ar[l,swap,"\varphi_{n-1}"]
&
P_{u_n}
\ar[l,swap,"\varphi_n"] .
\end{tikzcd}
\]

\begin{definition}[Intensional universal Mealy program machine]
The \textbf{intensional universal Mealy program machine} is the pullback of the extensional
universal evaluator along the decoder:
\[
        \mathbb P_{B,P}
        :=
        \operatorname{decode}_{B,P}^{\ast}\mathbb U_{B,P}.
\]
Thus
\[
        \mathbb P_{B,P}:\mathsf{Prog}_B(P)\rightsquigarrow B
\]
has carrier \(P\), and executes a formal program by decoding it to its compiled transition-output
operator.
\end{definition}

\[
\begin{tikzcd}[column sep=large,row sep=large]
\mathsf{Prog}_B(P)
\ar[r,"\operatorname{decode}_{B,P}"]
\ar[dr,swap,"\mathbb P_{B,P}"]
&
\mathsf U_B(P)
\ar[d,"\mathbb U_{B,P}"]
\\
&
B .
\end{tikzcd}
\]

\begin{theorem}[Straight-line compilation]
\label{thm:straight-line-compilation}
Every finite straight-line program in universal instructions has a canonical compiled universal
instruction.  Explicitly, for every word
\[
        w=
        [\varphi_1^{\operatorname{op}},\ldots,\varphi_n^{\operatorname{op}}]
\]
in \(\mathsf{Prog}_B(P)\),
\[
        \operatorname{decode}_{B,P}(w)
        =
        (\varphi_1\diamond\cdots\diamond\varphi_n)^{\operatorname{op}}.
\]
Executing \(w\) in the intensional universal machine has the same transition-output effect as
executing \(\operatorname{decode}_{B,P}(w)\) in the extensional universal machine:
\[
        \mathbb P_{B,P}(w)
        =
        \mathbb U_{B,P}\bigl(\operatorname{decode}_{B,P}(w)\bigr).
\]
Moreover, the compiled instruction is unique with this transition-output effect.
\end{theorem}

\begin{proof}
The free-category universal property gives a unique functor
\[
        \operatorname{decode}_{B,P}:
        \mathsf{Prog}_B(P)\to\mathsf U_B(P)
\]
extending the identity on primitive universal instructions.  Since
\(\operatorname{decode}_{B,P}\) is a functor, it sends a finite word to the categorical composite
of its letters in \(\mathsf U_B(P)\).

Composition in
\[
        \mathsf U_B(P)=\mathrm{KlEnd}_B(P)^{\operatorname{op}}
\]
is opposite to Kleisli composition in \(\mathrm{KlEnd}_B(P)\).  Therefore
\[
        [\varphi_1^{\operatorname{op}},\ldots,\varphi_n^{\operatorname{op}}]
\]
decodes to
\[
        (\varphi_1\diamond\cdots\diamond\varphi_n)^{\operatorname{op}},
\]
where the underlying Kleisli operator executes
\[
        \varphi_n,\ \varphi_{n-1},\ \ldots,\ \varphi_1
\]
in that order.

By definition,
\[
        \mathbb P_{B,P}
        =
        \operatorname{decode}_{B,P}^{\ast}\mathbb U_{B,P}.
\]
Therefore executing \(w\) in \(\mathbb P_{B,P}\) is executing its decoded instruction in
\(\mathbb U_{B,P}\).  If another universal instruction has the same transition-output effect as
\(\operatorname{decode}_{B,P}(w)\), then Theorem~\ref{thm:operational-full-abstraction} identifies
it with \(\operatorname{decode}_{B,P}(w)\).  Hence the compiled instruction is unique.
\end{proof}

Thus
\[
        \operatorname{decode}_{B,P}
\]
is the canonical straight-line compiler:
\[
        \text{formal program}
        \longmapsto
        \text{single extensional transition-output instruction}.
\]

\begin{corollary}[One-letter normal forms]
\label{cor:one-letter-normal-forms}
Every formal universal program has a unique one-letter normal form up to operational equality.
Specifically, for every
\[
        w\in \mathsf{Prog}_B(P)_1,
\]
the one-letter word
\[
        [\operatorname{decode}_{B,P}(w)]
\]
has the same evaluated transition-output effect as \(w\), and its letter is uniquely determined in
\(\mathsf U_B(P)\).
\end{corollary}

\begin{proof}
By Theorem~\ref{thm:straight-line-compilation}, \(w\) compiles to the instruction
\[
        \operatorname{decode}_{B,P}(w).
\]
The one-letter word
\[
        [\operatorname{decode}_{B,P}(w)]
\]
decodes to the same instruction.  Operational full abstraction gives uniqueness of the compiled
instruction.
\end{proof}

The corollary is the formal version of straight-line program compilation.  The intensional object
remembers the program as written.  The extensional object remembers only its compiled effect.

\section{The extensional quotient of the intensional program category}
\label{sec:extensional-quotient-intensional-language}

The decoder induces the operational quotient of the intensional program category.

\begin{definition}[Operational congruence on universal programs]
Let \(w,w'\) be arrows of \(\mathsf{Prog}_B(P)\) with the same source and target.  They are
\textbf{operationally congruent} if
\[
        \operatorname{decode}_{B,P}(w)
        =
        \operatorname{decode}_{B,P}(w')
\]
in \(\mathsf U_B(P)\).  We write
\[
        w\equiv_{\operatorname{op}} w'.
\]
\end{definition}

\begin{theorem}[Extensional quotient of the intensional program category]
\label{thm:prog-quotient-u}
The decoder
\[
        \operatorname{decode}_{B,P}:
        \mathsf{Prog}_B(P)\longrightarrow\mathsf U_B(P)
\]
exhibits \(\mathsf U_B(P)\) as the operational quotient of
\(\mathsf{Prog}_B(P)\):
\[
        \mathsf U_B(P)
        \cong
        \mathsf{Prog}_B(P)/{\equiv_{\operatorname{op}}}.
\]
\end{theorem}

\begin{proof}
The functor
\[
        \operatorname{decode}_{B,P}
\]
is identity-on-objects.  Let
\[
        j_1:\mathsf U_B(P)_1\longrightarrow \mathsf{Prog}_B(P)_1
\]
be the one-letter inclusion, sending an instruction \(\theta\) to the word \([\theta]\).  Then
\[
        \operatorname{decode}_{B,P,1}\,j_1
        =
        1_{\mathsf U_B(P)_1}.
\]
Hence the arrow map of the decoder is a split epimorphism, and therefore a regular epimorphism in
a regular category.

The kernel congruence of the arrow map of the decoder is precisely
\[
        w\equiv_{\operatorname{op}}w'.
\]
Since the decoder is a regular epimorphism on arrows and identity-on-objects, its regular image is
all of \(\mathsf U_B(P)\).  By the effective operational quotient condition in
Definition~\ref{def:universal-mealy-program-base}, the quotient by its kernel congruence is
isomorphic to \(\mathsf U_B(P)\), with quotient functor
\[
        \operatorname{decode}_{B,P}.
\]
\end{proof}

The quotient picture is:
\[
\begin{tikzcd}[column sep=large,row sep=large]
\mathsf{Prog}_B(P)
\ar[rr,"\operatorname{decode}_{B,P}"]
\ar[dr,swap,two heads]
&
&
\mathsf U_B(P)
\\
&
\mathsf{Prog}_B(P)/{\equiv_{\operatorname{op}}}
\ar[ur,swap,"\cong"]
&
.
\end{tikzcd}
\]

This theorem identifies the source of extensionality.  The universal operator category
\[
        \mathsf U_B(P)
\]
is not the primitive intensional language.  It is the fully abstract operational quotient of the
canonical intensional straight-line language
\[
        \mathsf{Prog}_B(P).
\]

\begin{corollary}[Full abstraction of the quotient]
\label{cor:full-abstraction-quotient}
The quotient
\[
        \mathsf{Prog}_B(P)\twoheadrightarrow\mathsf U_B(P)
\]
is fully abstract for carrier-level transition-output effect.  That is,
\[
        w\equiv_{\operatorname{op}}w'
\]
if and only if \(w\) and \(w'\) have the same evaluated transition-output effect on \(P\).
\end{corollary}

\begin{proof}
By definition,
\[
        w\equiv_{\operatorname{op}}w'
\]
means
\[
        \operatorname{decode}_{B,P}(w)
        =
        \operatorname{decode}_{B,P}(w').
\]
By Theorem~\ref{thm:operational-full-abstraction}, equality of decoded universal instructions is
equivalent to equality of their evaluated transition-output effects.  Since
\[
        \mathbb P_{B,P}
        =
        \operatorname{decode}_{B,P}^{\ast}\mathbb U_{B,P},
\]
the evaluated effect of a formal program is the evaluated effect of its decoded instruction.
\end{proof}

\section{Generated intensional languages}
\label{sec:generated-intensional-languages}

The canonical intensional program category
\[
        \mathsf{Prog}_B(P)
\]
has all universal operators as primitive instructions.  It is canonical, but a concrete programming
language is obtained by choosing a graph of primitive instruction symbols and interpreting those
symbols in the universal operator category.

Let
\[
        G\rightrightarrows I
\]
be an internal graph.  A generalized arrow of \(G\) is a primitive instruction symbol with a source
and target object of \(I\).

\begin{definition}[Primitive decoding]
A \textbf{primitive decoding} of \(G\) into the universal Mealy evaluator is a map of internal
graphs
\[
        \gamma:G\longrightarrow \mathsf U_B(P)_{\mathrm{gr}}.
\]
Equivalently, it assigns to each primitive instruction symbol an extensional transition-output
operator on \(P\), with matching source and target types.
\end{definition}

Let
\[
        \mathrm{FreeCat}(G)
\]
be the free internal category on \(G\).

\begin{theorem}[Free intensional extension]
\label{thm:free-intensional-extension}
Every primitive decoding
\[
        \gamma:G\longrightarrow \mathsf U_B(P)_{\mathrm{gr}}
\]
extends uniquely to an identity-on-objects internal functor
\[
        \widehat{\gamma}:
        \mathrm{FreeCat}(G)
        \longrightarrow
        \mathsf U_B(P).
\]
The induced Mealy morphism
\[
        \widehat{\gamma}^{\ast}\mathbb U_{B,P}
        :
        \mathrm{FreeCat}(G)\rightsquigarrow B
\]
is the Mealy machine whose programs are formal words in the primitive instructions of \(G\), and
whose execution is obtained by universal evaluation after decoding.
\end{theorem}

\begin{proof}
The primitive decoding
\[
        \gamma:G\to\mathsf U_B(P)_{\mathrm{gr}}
\]
is a graph map into the underlying graph of the internal category \(\mathsf U_B(P)\).  By the
universal property of \(\mathrm{FreeCat}(G)\), it extends uniquely to an internal functor
\[
        \widehat{\gamma}:\mathrm{FreeCat}(G)\to\mathsf U_B(P).
\]
Restricting the universal evaluator
\[
        \mathbb U_{B,P}:\mathsf U_B(P)\rightsquigarrow B
\]
along \(\widehat{\gamma}\) gives a Mealy morphism
\[
        \widehat{\gamma}^{\ast}\mathbb U_{B,P}
        :
        \mathrm{FreeCat}(G)\rightsquigarrow B.
\]
Its carrier is \(P\), its input programs are formal words in \(G\), and its transition-output
operation is obtained by decoding those words into \(\mathsf U_B(P)\).
\end{proof}

Thus a programming language over \(P\) is not another universal operator category.  It is an
intensional graph or category lying over the universal operator category:
\[
\begin{tikzcd}[column sep=large,row sep=large]
G
\ar[r,"\gamma"]
\ar[d]
&
\mathsf U_B(P)_{\mathrm{gr}}
\ar[d]
\\
\mathrm{FreeCat}(G)
\ar[r,"\widehat{\gamma}"]
&
\mathsf U_B(P).
\end{tikzcd}
\]

The induced evaluator fits into:
\[
\begin{tikzcd}[column sep=large,row sep=large]
\mathrm{FreeCat}(G)
\ar[r,"\widehat{\gamma}"]
\ar[dr,swap,"\widehat{\gamma}^{\ast}\mathbb U_{B,P}"]
&
\mathsf U_B(P)
\ar[d,"\mathbb U_{B,P}"]
\\
&
B .
\end{tikzcd}
\]

The extensional operational image of such a language is
\[
        \operatorname{RegIm}_I(\widehat{\gamma})
        \hookrightarrow
        \mathsf U_B(P).
\]
This is the generated operator fragment of the language.

\begin{definition}[Generated operational fragment]
Let
\[
        \gamma:G\to\mathsf U_B(P)_{\mathrm{gr}}
\]
be a primitive decoding.  The \textbf{generated operational fragment} of \(\gamma\) is
\[
        \operatorname{GenOp}(\gamma)
        :=
        \operatorname{RegIm}_I
        \bigl(
            \widehat{\gamma}:\mathrm{FreeCat}(G)\to\mathsf U_B(P)
        \bigr).
\]
\end{definition}

\[
\begin{tikzcd}[column sep=large,row sep=large]
\mathrm{FreeCat}(G)
\ar[rr,"\widehat{\gamma}"]
\ar[dr,swap,two heads]
&
&
\mathsf U_B(P)
\\
&
\operatorname{GenOp}(\gamma)
\ar[ur,swap,hook]
&
.
\end{tikzcd}
\]

\begin{proposition}[Image universal property of generated operational fragments]
\label{prop:minimality-generated-operational-fragment}
The subcategory
\[
        \operatorname{GenOp}(\gamma)\hookrightarrow \mathsf U_B(P)
\]
is the identity-on-objects regular image of
\[
        \widehat{\gamma}:\mathrm{FreeCat}(G)\to\mathsf U_B(P).
\]
Consequently, every identity-on-objects internal subcategory
\[
        \mathcal D\hookrightarrow \mathsf U_B(P)
\]
through which \(\widehat{\gamma}\) factors contains \(\operatorname{GenOp}(\gamma)\).
\end{proposition}

\begin{proof}
The first statement is the definition of \(\operatorname{GenOp}(\gamma)\).  Since
\(\widehat{\gamma}\) is an internal functor, its regular image is an identity-on-objects internal
subcategory of \(\mathsf U_B(P)\).  If
\[
        \widehat{\gamma}:\mathrm{FreeCat}(G)\to\mathcal D\hookrightarrow \mathsf U_B(P)
\]
is a factorization through an identity-on-objects subcategory \(\mathcal D\), then the regular
image of \(\widehat{\gamma}\) in \(\mathsf U_B(P)\) factors through \(\mathcal D\).  Hence
\(\mathcal D\) contains \(\operatorname{GenOp}(\gamma)\).
\end{proof}

This proposition is the categorical form of the distinction between a primitive instruction set
and the operator category generated by that instruction set.

\section{Carrier images, syntactic images, and behaviour}
\label{sec:carrier-syntactic-behaviour-images}

The preceding section separates syntax from operators.  This section relates that separation to
the image constructions used earlier in the manuscript.

Let
\[
        d:C\longrightarrow\mathsf U_B(P)
\]
be an identity-on-objects functor from an intensional instruction category \(C\).  The
\textbf{operational image} of \(C\) on \(P\) is
\[
        \operatorname{OpIm}_P(C)
        :=
        \operatorname{RegIm}_I(d)
        \hookrightarrow
        \mathsf U_B(P).
\]
This is the extensional fragment generated by the intensional language \(C\).

For a Mealy morphism
\[
        \Phi:A\rightsquigarrow B
\]
with carrier \(P\), the classifying functor
\[
        \ulcorner\Phi\urcorner:A\to\mathsf U_B(P)
\]
has operational image
\[
        \operatorname{OpIm}_P(A)
        =
        \operatorname{CarOp}(\Phi)^{\operatorname{op}}.
\]

For a specification \(K\), when the carrier is the derivative-closed behaviour carrier
\[
        \operatorname{Der}_0(K),
\]
the same construction gives the syntactic transition-output category:
\[
        \mathrm{SynTO}_{A,B}(K)^{\operatorname{op}}
        \hookrightarrow
        \mathsf U_B(\operatorname{Der}_0(K)).
\]

The distinction is:
\[
\begin{array}{c|c|c}
\text{object} & \text{formed from} & \text{role} \\ \hline
A,\ C,\ \mathrm{FreeCat}(G)
    & \text{syntax or chosen primitives}
    & \text{intensional program category} \\
\operatorname{CarOp}(\Phi)
    & \text{carrier action of }\Phi
    & \text{operators generated on }P \\
\mathrm{SynTO}_{A,B}(K)
    & \text{action on }\operatorname{Der}_0(K)
    & \text{syntactic operator image of a specification} \\
\mathrm{KlEnd}_B(P)
    & \text{all fibrewise Kleisli endomorphisms of }P
    & \text{full extensional operator category.}
\end{array}
\]

There is also a behavioural level.  Terminal residual semantics sends a state of a Mealy morphism
to its residual behaviour.  Thus two implementation-level operators may become behaviourally
identified after passing from carrier states to residual behaviours.  The universal operator
category lives before this black-boxing step:
\[
\begin{tikzcd}[column sep=large,row sep=large]
\text{syntax}
\ar[r]
&
\text{carrier operators}
\ar[r]
&
\text{residual behaviours}
\\
A
\ar[r,"\ulcorner\Phi\urcorner"]
&
\mathsf U_B(P)
\ar[r,dashed]
&
\operatorname{Op}_{\mathrm{beh}}(\Phi)^{\operatorname{op}} .
\end{tikzcd}
\]
The dashed arrow denotes the behavioural quotient induced by residual semantics on the operators
generated by the given input action.  It is not a functor out of the full universal operator
category in general.

\begin{definition}[Behavioural transition-output image]
Let
\[
        \Phi:A\rightsquigarrow B
\]
have carrier \(P\), and let
\[
        \beta_P:P\longrightarrow \Beh^{\operatorname{can}}_{A,B}
\]
be its residual semantics map.  Let
\[
        \Beh(\Phi)\hookrightarrow \Beh^{\operatorname{can}}_{A,B}
\]
be the regular image of \(\beta_P\).  Since \(\beta_P\) is a strict semantic homomorphism, the
carrier-level transition-output action descends to the behavioural image.  The
\textbf{behavioural transition-output image} of \(\Phi\) is
\[
        \operatorname{Op}_{\mathrm{beh}}(\Phi)
        :=
        \operatorname{RegIm}_I
        \left(
            A^{\operatorname{op}}
            \longrightarrow
            \mathrm{KlEnd}_B(\Beh(\Phi))
        \right).
\]
\end{definition}

\begin{proposition}[Behavioural collapse of carrier operators]
\label{prop:behavioural-collapse-carrier-operators}
There is a canonical identity-on-objects comparison
\[
        \operatorname{CarOp}(\Phi)
        \longrightarrow
        \operatorname{Op}_{\mathrm{beh}}(\Phi).
\]
When the induced action on the behavioural image is the regular quotient of the carrier action,
this comparison is a regular epimorphism.  Its kernel identifies precisely those generated carrier
operators whose induced effects agree on the residual behavioural image
\[
        \Beh(\Phi).
\]
\end{proposition}

\begin{proof}
The carrier operational image is generated by the action
\[
        A^{\operatorname{op}}\to\mathrm{KlEnd}_B(P).
\]
The residual semantics map
\[
        \beta_P:P\to\Beh^{\operatorname{can}}_{A,B}
\]
is a strict semantic homomorphism, so it is compatible with the transition-output action of
\(A^{\operatorname{op}}\).  Taking the regular image
\[
        \Beh(\Phi)\hookrightarrow\Beh^{\operatorname{can}}_{A,B}
\]
and restricting the induced action gives a transition-output action
\[
        A^{\operatorname{op}}\to\mathrm{KlEnd}_B(\Beh(\Phi)).
\]

We now check that the carrier operational quotient maps to the behavioural operational quotient.
Suppose \(a,b:u\to v\) have the same image in \(\operatorname{CarOp}(\Phi)\).  Equivalently, their
induced Kleisli operators on \(P\) are equal.  Since \(\beta_P\) is strict, applying \(\beta_P\)
after either transition gives the same transition on the residual image.  Pulling back along the
regular-epimorphic cover
\[
        P\twoheadrightarrow \Beh(\Phi)
\]
therefore makes the two induced behavioural operators equal.  By
Lemma~\ref{lem:regular-epi-detection-operator-equality}, the induced operators on
\(\Beh(\Phi)\) are equal.  Hence the kernel of the carrier action is contained in the kernel of
the behavioural-image action.

It follows, by the universal property of the carrier operational image, that the behavioural-image
action factors through a canonical identity-on-objects comparison
\[
        \operatorname{CarOp}(\Phi)
        \longrightarrow
        \operatorname{Op}_{\mathrm{beh}}(\Phi).
\]
If the behavioural-image action is the regular quotient of the carrier action, then the induced
comparison of regular images is a regular epimorphism.  Its kernel is exactly the relation of
agreeing on the residual behavioural image.
\end{proof}

This proposition records the distinction between implementation-level operational equality and
black-box behavioural equality.  The universal category \(\mathsf U_B(P)\) is extensional at the
operator level, not necessarily at the behavioural level.

\section{Finite Mealy--Cayley representation}
\label{sec:finite-mealy-cayley-representation}

The one-object finite-state case gives a concrete theorem showing that the abstract construction
recovers the full transformation-output monoid.

Let \(N\) be a monoid, regarded as a one-object category, and let \(Q\) be a finite set.  Set
\[
        I=1,
        \qquad
        B=N,
        \qquad
        P=Q.
\]
Then a fibrewise output-comparison Kleisli endomorphism of \(Q\) consists of a function
\[
        f:Q\to Q
\]
together with an output labelling
\[
        \omega:Q\to N.
\]

We write elements of
\[
        N^Q\rtimes\operatorname{End}(Q)
\]
as pairs
\[
        (f,\omega),
\]
with the endomap displayed first.

Write
\[
        (f,\omega)\star(g,\nu)
\]
for the operation which executes \((f,\omega)\) first and \((g,\nu)\) second.

\begin{theorem}[Finite Mealy--Cayley representation]
\label{thm:finite-mealy-cayley}
There is an isomorphism of monoids
\[
        \mathrm{KlEnd}_N(Q)
        \cong
        N^Q\rtimes \operatorname{End}(Q),
\]
where the semidirect product multiplication is
\[
        (f,\omega)\star(g,\nu)
        =
        \left(g\circ f,\;q\mapsto \omega(q)\circ\nu(fq)\right)
\]
in ordinary categorical composition notation.

Equivalently, using the manuscript's multiplicative convention \(b\circ a=ba\),
\[
        (f,\omega)\star(g,\nu)
        =
        \left(g\circ f,\;q\mapsto \omega(q)\,\nu(fq)\right).
\]
The unit is
\[
        (1_Q,q\mapsto 1_N).
\]
Consequently, every deterministic \(N\)-output Mealy action of an input monoid \(M\) on \(Q\) is
equivalently a monoid homomorphism
\[
        M^{\operatorname{op}}
        \longrightarrow
        N^Q\rtimes\operatorname{End}(Q).
\]
\end{theorem}

\begin{proof}
Since \(I=1\), there is only one fibre of the carrier, namely \(Q\).  Since \(B=N\) is a
one-object category, the output-comparison arrow associated to a transition from \(q\) to \(f(q)\)
is an element of \(N\).  Hence a Kleisli endomorphism is precisely a pair
\[
        (f,\omega),
        \qquad
        f:Q\to Q,
        \quad
        \omega:Q\to N.
\]

Executing \((f,\omega)\) first and \((g,\nu)\) second sends
\[
        q
        \longmapsto
        f(q)
        \longmapsto
        g(f(q)).
\]
Thus the state component of the composite is
\[
        g\circ f.
\]
The output component first contributes
\[
        \omega(q)
\]
and then contributes
\[
        \nu(f(q)).
\]
With ordinary categorical composition, the resulting output comparison is
\[
        \omega(q)\circ\nu(f(q)).
\]
Using the manuscript's multiplicative convention \(b\circ a=ba\), this is
\[
        \omega(q)\,\nu(fq).
\]
The identity operation is
\[
        (1_Q,q\mapsto 1_N).
\]
Therefore
\[
        \mathrm{KlEnd}_N(Q)
        \cong
        N^Q\rtimes\operatorname{End}(Q).
\]

A deterministic \(N\)-output Mealy action of an input monoid \(M\) on \(Q\) is a one-object Mealy
structure, hence a Kleisli action
\[
        M^{\operatorname{op}}\to\mathrm{KlEnd}_N(Q).
\]
Using the displayed isomorphism gives the required homomorphism
\[
        M^{\operatorname{op}}\to N^Q\rtimes\operatorname{End}(Q).
\]
\end{proof}

This theorem is the Mealy analogue of the Cayley representation by transformations.  The ordinary
full transformation monoid
\[
        \operatorname{End}(Q)
\]
is replaced by the full transformation-output monoid
\[
        N^Q\rtimes\operatorname{End}(Q).
\]

The corresponding universal finite-state Mealy morphism is
\[
        \mathbb U_{N,Q}:
        \bigl(N^Q\rtimes \operatorname{End}(Q)\bigr)^{\operatorname{op}}
        \rightsquigarrow N.
\]
It executes
\[
        \delta_{\mathbb U}((f,\omega)^{\operatorname{op}},q)=f(q),
        \qquad
        \omega_{\mathbb U}((f,\omega)^{\operatorname{op}},q)=\omega(q).
\]

The intensional universal program monoid is
\[
        \mathrm{FreeMon}\bigl(N^Q\rtimes\operatorname{End}(Q)\bigr).
\]
Its decoder sends a word
\[
        [(f_1,\omega_1)^{\operatorname{op}},\ldots,(f_n,\omega_n)^{\operatorname{op}}]
\]
to
\[
        \bigl(
        (f_n,\omega_n)\star\cdots\star(f_1,\omega_1)
        \bigr)^{\operatorname{op}}.
\]
This is the same order as in the general formula: the last letter of the input word is executed
first at the carrier level.

Thus the finite-state case has the same three levels:
\[
\begin{array}{c|c}
\text{level} & \text{finite one-object form} \\ \hline
\text{intensional programs}
    & \mathrm{FreeMon}\bigl(N^Q\rtimes\operatorname{End}(Q)\bigr) \\
\text{extensional operators}
    & N^Q\rtimes\operatorname{End}(Q) \\
\text{execution}
    & \mathbb U_{N,Q}:
        \bigl(N^Q\rtimes\operatorname{End}(Q)\bigr)^{\operatorname{op}}
        \rightsquigarrow N .
\end{array}
\]

\section{Lawvere diagonal obstruction to total self-coding}
\label{sec:lawvere-diagonal-obstruction}

The extensional universal operator category is canonical, but it is not a self-coded
stored-program universal computer.  In total settings, full extensional self-coding is obstructed
by diagonalization.

\begin{theorem}[Lawvere diagonal obstruction]
\label{thm:lawvere-diagonal-obstruction}
Let \(\mathcal E=\mathrm{Set}\), let \(P\) be a set, and suppose \(P\) has a fixed-point-free
endomap
\[
        s:P\longrightarrow P.
\]
Then there is no map
\[
        e:P\times P\longrightarrow P
\]
which is point-surjective for all endomaps \(P\to P\).  That is, there is no \(e\) such that for
every function
\[
        f:P\longrightarrow P
\]
there exists \(p\in P\) with
\[
        f(x)=e(p,x)
        \quad\text{for all }x\in P.
\]
\end{theorem}

\begin{proof}
Assume, for contradiction, that such an evaluator
\[
        e:P\times P\to P
\]
exists.  Define the diagonal function
\[
        d:P\to P,
        \qquad
        d(x)=s(e(x,x)).
\]
By point-surjectivity of \(e\), there exists \(p\in P\) such that
\[
        d(x)=e(p,x)
\]
for all \(x\).  Evaluating at \(x=p\), one obtains
\[
        e(p,p)=d(p)=s(e(p,p)).
\]
Thus \(e(p,p)\) is a fixed point of \(s\), contradicting the assumption that \(s\) is fixed-point
free.
\end{proof}

This is the Cantor--Lawvere phenomenon in computational form \cite{Lawvere1969}.  It explains why
the full total extensional operator object
\[
        \mathsf U_1(P)=\operatorname{End}(P)^{\operatorname{op}}
\]
cannot generally be coded by \(P\) itself.

In the one-object total case,
\[
        I=1,
        \qquad
        B=1,
\]
the extensional universal evaluator is ordinary evaluation
\[
        \operatorname{End}(P)\times P\longrightarrow P.
\]
A self-coded total evaluator would require a coding map
\[
        P\longrightarrow \operatorname{End}(P)
\]
which is surjective enough to represent all total endomaps.  The theorem says that this is
impossible whenever \(P\) has a fixed-point-free endomap.

For example, with
\[
        P=\mathbb N
\]
and
\[
        s(n)=n+1,
\]
there is no total self-indexed evaluator
\[
        e:\mathbb N\times\mathbb N\to\mathbb N
\]
representing all total functions \(\mathbb N\to\mathbb N\).

This does not contradict the existence of
\[
        \mathsf U_1(\mathbb N)=\operatorname{End}(\mathbb N)^{\operatorname{op}}.
\]
It says that the full extensional instruction object
\[
        \operatorname{End}(\mathbb N)
\]
cannot be replaced by the state object \(\mathbb N\) itself in a total extensional way.

Thus:
\[
\boxed{
\text{full extensional operator universality}
\neq
\text{total self-coded universal computation.}
}
\]
Classical self-coded universality requires intensional codes and, for ordinary recursion theory,
partial evaluation.

\begin{corollary}[No total primitive-recursive universal evaluator for primitive recursion]
\label{cor:no-total-pr-self-universal}
There is no primitive-recursive function
\[
        U:\mathbb N\times\mathbb N\to\mathbb N
\]
such that every unary primitive-recursive function
\[
        f:\mathbb N\to\mathbb N
\]
has an index \(e\in\mathbb N\) with
\[
        f(x)=U(e,x)
        \quad\text{for all }x.
\]
\end{corollary}

\begin{proof}
Suppose, for contradiction, that such a primitive-recursive evaluator \(U\) exists.  Define
\[
        g(x)=U(x,x)+1.
\]
Since \(U\) is primitive recursive, \(g\) is primitive recursive.  By universality, there exists an
index \(e_0\) such that
\[
        g(x)=U(e_0,x)
\]
for all \(x\).  Evaluating at \(x=e_0\), one obtains
\[
        U(e_0,e_0)+1
        =
        g(e_0)
        =
        U(e_0,e_0),
\]
a contradiction.
\end{proof}

The point is not that primitive recursion is irrelevant.  The point is that a total evaluator
cannot be both primitive recursive and universal for all unary primitive-recursive functions.
Classical universality escapes this diagonal obstruction through partiality, intensional coding, or
by using an evaluator outside the represented total class.

\section{Partiality, coding, and classical universality}
\label{sec:partiality-coding-classical-universality}

The constructions above are total.  They isolate the categorical universal operator evaluator and
its canonical intensional straight-line language.

Classical universal computation adds further structure:
\[
\begin{array}{c}
\text{intensional codes} \\
{}+\ \text{partial evaluation} \\
{}+\ \text{composition tracking} \\
{}+\ \text{parameterization} \\
{}+\ \text{finite or unbounded iteration.}
\end{array}
\]

\begin{definition}[Stable fibrewise partiality structure]
A \textbf{stable fibrewise partiality structure} on the slices of \(\mathcal E\) consists of
monads
\[
        L_X:\mathcal E/X\longrightarrow \mathcal E/X
\]
for the slices used below, stable under pullback in \(X\), together with Kleisli composition and
units compatible with reindexing.
\end{definition}

Given such a structure, a partial transition-output operator over \(I\times I\) is a Kleisli
section in the slice over \(P_{\mathrm{src}}\):
\[
        P_{\mathrm{src}}
        \longrightarrow
        L_{P_{\mathrm{src}}}\bigl(E_B(P)^{[2]}\bigr).
\]

\begin{definition}[Partial operator category]
A stable fibrewise partiality structure \(L\) is \textbf{admissible for the carrier \(P\) over
\(B\)} if the partial transition-output operators defined above admit identities and Kleisli
composition internally, forming an internal category
\[
        \mathrm{KlEnd}^L_B(P).
\]
The corresponding partial universal Mealy category is
\[
        \mathsf U^L_B(P)
        :=
        \mathrm{KlEnd}^L_B(P)^{\operatorname{op}}.
\]
\end{definition}

\begin{proposition}[Partial universal Mealy category]
\label{prop:partial-universal-mealy-category}
Let \(L\) be admissible for \(P\) over \(B\).  Then
\[
        \mathsf U^L_B(P):=\mathrm{KlEnd}^L_B(P)^{\operatorname{op}}
\]
represents partial Mealy structures on the fixed carrier \(P\).
\end{proposition}

\begin{proof}
A partial Mealy structure on \(P\) is a fibrewise partial Kleisli action
\[
        A^{\operatorname{op}}\longrightarrow \mathrm{KlEnd}^L_B(P).
\]
Taking opposites converts this datum into an identity-on-objects internal functor
\[
        A\longrightarrow
        \mathrm{KlEnd}^L_B(P)^{\operatorname{op}}
        =
        \mathsf U^L_B(P).
\]
The converse construction is obtained by opposing such a functor.  This is the partial analogue of
Theorem~\ref{thm:genericity-universal-mealy}.
\end{proof}

An effective or computable programming language is not obtained canonically from
\(\mathsf U^L_B(P)\) alone.  One must choose an effective operator subcategory
\[
        \mathcal C\hookrightarrow \mathrm{KlEnd}^L_B(P)
\]
and a code object mapping to its arrows.

\begin{definition}[Turing-style Mealy coding refinement]
A \textbf{Turing-style Mealy coding refinement} consists of:
\begin{enumerate}
    \item a stable fibrewise partiality structure \(L\);
    \item an internal category \(B\);
    \item a carrier
    \[
            I\xleftarrow{p}P\xrightarrow{q}B_0;
    \]
    \item an identity-on-objects effective operator subcategory
    \[
            \mathcal C\hookrightarrow \mathrm{KlEnd}^L_B(P);
    \]
    \item an object of typed codes
    \[
            c:\operatorname{Code}\to I\times I;
    \]
    \item a coding map over \(I\times I\)
    \[
            \gamma:\operatorname{Code}\to \mathcal C_1;
    \]
    \item a code-indexed one-step partial evaluator, defined by pulling back universal partial
    evaluation along \(\gamma\);
    \item identity and composition tracking on codes, meaning maps on codes whose images under
    \(\gamma\) are the identity and Kleisli-composite operators in \(\mathcal C\).
\end{enumerate}
It is \textbf{universal for \(\mathcal C\)} when
\[
        \gamma:\operatorname{Code}\to\mathcal C_1
\]
is a regular epimorphism in \(\mathcal E/(I\times I)\), or a split epimorphism when global chosen
codes are required.
\end{definition}

\begin{proposition}[Coded partial evaluation]
\label{prop:coded-partial-evaluation}
A Turing-style Mealy coding refinement determines a partial Mealy morphism
\[
        \mathcal C^{\operatorname{op}}\rightsquigarrow B
\]
whose primitive instructions are the operators of \(\mathcal C\).  The code object locally names
these instructions, and the code-indexed one-step evaluator is the pullback of universal partial
evaluation along
\[
        \gamma:\operatorname{Code}\to\mathcal C_1.
\]
With either a free category of coded programs or explicit identity and composition tracking on
codes, this extends to a coded partial program evaluator.
\end{proposition}

\begin{proof}
The inclusion
\[
        \mathcal C\hookrightarrow \mathrm{KlEnd}^L_B(P)
\]
opposes to an inclusion
\[
        \mathcal C^{\operatorname{op}}\hookrightarrow \mathsf U^L_B(P).
\]
Restricting the partial universal Mealy evaluator along this inclusion gives a partial Mealy
morphism
\[
        \mathcal C^{\operatorname{op}}\rightsquigarrow B.
\]
The coding map
\[
        \gamma:\operatorname{Code}\to\mathcal C_1
\]
locally names the arrows of \(\mathcal C\).  Pulling back universal partial evaluation along
\(\gamma\) gives the code-indexed one-step evaluator.  If the codes freely generate a coded program
category, or if the code object carries identity and composition tracking, then the one-step
evaluator extends to a coded partial program evaluator.
\end{proof}

This situation is summarized by:
\[
\begin{tikzcd}[column sep=large,row sep=large]
\operatorname{Code}
\ar[r,"\gamma"]
\ar[d]
&
\mathcal C_1
\ar[d,hook]
\\
I\times I
&
\mathrm{KlEnd}^L_B(P)_1
\ar[l]
\end{tikzcd}
\]

In the one-object case
\[
        I=1,
        \qquad
        P=\mathbb N,
        \qquad
        B=1,
\]
choosing the effective subcategory of partial recursive endomaps and an acceptable enumeration
\[
        \gamma:\mathbb N\longrightarrow \mathcal C_1
\]
recovers the ordinary universal partial recursive evaluator
\[
        (e,x)\longmapsto \varphi_e(x)
\]
as a coded pullback of partial universal Mealy evaluation \cite{Kleene1952,Rogers1967}.  The
\(s\)-\(m\)-\(n\) theorem supplies parameterization, and computable manipulation of indices tracks
composition.

Categorical accounts of such partial-computability structures include Turing categories and
related abstract-computability frameworks \cite{CockettHofstra2008}.  These structures refine the
total construction
\[
        \mathsf{Prog}_B(P)\to\mathsf U_B(P)
\]
by choosing codes and partial effective operator fragments.

\section{Examples}
\label{sec:universal-mealy-examples}

\subsection{Total functions in \(\mathrm{Set}\)}

Let
\[
        \mathcal E=\mathrm{Set},
        \qquad
        I=1,
        \qquad
        B=1,
        \qquad
        P=X.
\]
Then
\[
        \mathrm{KlEnd}_1(X)\cong \operatorname{End}(X)
\]
as a one-object category.  The extensional universal Mealy category is
\[
        \mathsf U_1(X)=\operatorname{End}(X)^{\operatorname{op}}.
\]
The universal evaluator is ordinary evaluation:
\[
        \delta_{\mathbb U}(f^{\operatorname{op}},x)=f(x).
\]

The intensional universal program category is the free monoid on the set of endomaps:
\[
        \mathsf{Prog}_1(X)
        \cong
        \mathrm{FreeMon}(\operatorname{End}(X)).
\]
A program is a finite word
\[
        [f_1^{\operatorname{op}},\ldots,f_n^{\operatorname{op}}].
\]
The decoder is ordinary composition with the opposite-category order:
\[
        [f_1^{\operatorname{op}},\ldots,f_n^{\operatorname{op}}]
        \longmapsto
        (f_1\circ\cdots\circ f_n)^{\operatorname{op}},
\]
so the carrier-level execution applies
\[
        f_n,\ f_{n-1},\ \ldots,\ f_1
\]
in that order.

Thus
\[
        \operatorname{End}(X)
\]
is the extensional quotient of the free straight-line language on all primitive transformations of
\(X\).

For \(X=\mathbb N\), this is not an effective computability object.  It contains all functions
\[
        \mathbb N\to\mathbb N,
\]
including noncomputable functions.  The construction is a total extensional operator universe, not
an effective programming language.

\subsection{Finite-state transformation-output monoids}

Let
\[
        I=1,
\]
let \(Q\) be a finite set, and let \(B\) be the one-object category associated to a monoid \(N\).
By Theorem~\ref{thm:finite-mealy-cayley},
\[
        \mathrm{KlEnd}_N(Q)
        \cong
        N^Q\rtimes \operatorname{End}(Q).
\]
The extensional universal finite-state Mealy morphism is
\[
        \mathbb U_{N,Q}:
        \bigl(N^Q\rtimes \operatorname{End}(Q)\bigr)^{\operatorname{op}}
        \rightsquigarrow N.
\]
The intensional universal program monoid is
\[
        \mathrm{FreeMon}\bigl(N^Q\rtimes\operatorname{End}(Q)\bigr).
\]
The decoder sends a word to the corresponding semidirect product composite, in the reversed
\(\star\)-order specified in Section~\ref{sec:finite-mealy-cayley-representation}.  This is
finite-state operational universality.  It is not Turing universality, since the carrier is finite.

\subsection{Generated instruction languages}

Let
\[
        G\rightrightarrows I
\]
be a chosen graph of primitive instructions, and let
\[
        \gamma:G\to\mathsf U_B(P)_{\mathrm{gr}}
\]
assign to each primitive instruction a transition-output operator.

Then
\[
        \mathrm{FreeCat}(G)
\]
is the intensional program category generated by those primitives.  Its decoder
\[
        \widehat{\gamma}:\mathrm{FreeCat}(G)\to\mathsf U_B(P)
\]
is the unique extension of \(\gamma\).  Its operational image
\[
        \operatorname{RegIm}_I(\widehat{\gamma})
        \hookrightarrow
        \mathsf U_B(P)
\]
is the extensional operator category generated by those primitives.

This is the general pattern for concrete programming languages, instruction sets, gate sets,
rewrite systems, register-machine primitives, and circuit bases.

\subsection{Presheaf and sheaf bases}

Let
\[
        \mathcal E=\mathrm{Set}^{\mathcal C^{\operatorname{op}}}
\]
be a presheaf topos.  Since \(\mathcal E\) is locally cartesian closed,
\[
        \mathrm{KlEnd}_B(P)
\]
exists internally for every internal category \(B\) and carrier \(P\).

The construction is contextual rather than merely pointwise.  At a stage \(c\), generalized
elements of
\[
        \mathrm{KlEnd}_B(P)_1
\]
are compatible families of fibrewise transition-output operators over all restrictions into \(c\).
Thus presheaf and sheaf bases give context-indexed universal operator categories.

The corresponding intensional program category
\[
        \mathsf{Prog}_B(P)
\]
is the internal free category on the contextual graph of universal instructions.  It records formal
contextual programs before they are compiled to contextual transition-output operators.

\subsection{Self-coding obstruction}

In the total one-object case,
\[
        \mathsf U_1(P)=\operatorname{End}(P)^{\operatorname{op}}.
\]
If \(P\) has a fixed-point-free endomap, then
Theorem~\ref{thm:lawvere-diagonal-obstruction} prevents \(\operatorname{End}(P)\) from being
represented by \(P\) itself in a total extensional way.

Thus a stored-program universal computer is not obtained by setting
\[
        \operatorname{Code}=P
\]
and demanding a total enumeration of all extensional endomaps.  One must instead use an
intensional code object and, for classical computation, partial evaluation.

\section{Non-examples and obstructions}
\label{sec:nonexamples-obstructions-universal-mealy}

\subsection{Finitely complete bases without operator objects}

Finite limits are sufficient to define spans, internal categories, and ordinary Mealy morphisms.
They are not sufficient to construct
\[
        \mathrm{KlEnd}_B(P)
\]
in general.  The latter requires dependent products representing fibrewise operator families.  In
a finitely complete category where the relevant source-state maps
\[
        \pi_{\mathrm{src}}:P_{\mathrm{src}}\to I\times I
\]
are not exponentiable in the required slices, the extensional universal Mealy category
\[
        \mathsf U_B(P)
\]
is not defined by the construction of this chapter.

\subsection{The extensional universal category is not syntax}

The category
\[
        \mathsf U_B(P)
\]
is not a universal programming language in the intensional sense.  Its arrows are already
operators.  Formal programs live instead in
\[
        \mathsf{Prog}_B(P)
        =
        \mathrm{FreeCat}\bigl(\mathsf U_B(P)_{\mathrm{gr}}\bigr),
\]
or in a generated language
\[
        \mathrm{FreeCat}(G)
\]
lying over \(\mathsf U_B(P)\).

Thus the correct reading is:
\[
        \mathsf{Prog}_B(P)
        \longrightarrow
        \mathsf U_B(P)
        \longrightarrow
        B.
\]
The first map compiles intensional programs.  The second Mealy morphism evaluates extensional
operators.

\subsection{The full \(\mathrm{Set}\)-endomorphism object}

The object
\[
        \mathrm{KlEnd}_1(\mathbb N)
        =
        \mathbb N^{\mathbb N}
\]
in \(\mathrm{Set}\) is not an effective computability object.  It is a correct total extensional
operator object, but it contains noncomputable functions.  Therefore full operator universality is
not the same as computability.

\subsection{No total self-coding of full extensional universality}

The Lawvere diagonal obstruction shows that, when \(P\) has a fixed-point-free endomap, there is no
total self-indexed evaluator
\[
        P\times P\to P
\]
which represents all total endomaps \(P\to P\).  Thus the full extensional universal evaluator
cannot generally be made into a total stored-program evaluator with codes equal to states.

\subsection{Partiality without coding}

A lift monad or partial-map structure gives partial maps, but does not by itself give classical
universal computation.  One also needs:
\[
\begin{array}{ll}
1. & \text{an object of intensional codes;}\\
2. & \text{an evaluator induced by those codes;}\\
3. & \text{closure of codes under composition;}\\
4. & \text{parameterization or substitution of data into codes;}\\
5. & \text{finite path or trace closure for repeated execution;}\\
6. & \text{partial or fixed-point structure for divergence and unbounded search.}
\end{array}
\]
Partiality is therefore necessary for comparison with classical recursion theory, but it is not
sufficient by itself.

\section{Chapter summary}
\label{sec:universal-mealy-summary}

This chapter separated intensional programs from extensional transition-output operators.

The extensional universal Mealy category of a carrier
\[
        I\xleftarrow{p}P\xrightarrow{q}B_0
\]
is
\[
        \mathsf U_B(P)
        =
        \mathrm{KlEnd}_B(P)^{\operatorname{op}}.
\]
Its associated universal evaluator is
\[
        \mathbb U_{B,P}:\mathsf U_B(P)\rightsquigarrow B.
\]
It executes transition-output operators by evaluation.

The genericity theorem says that every Mealy morphism
\[
        \Phi:A\rightsquigarrow B
\]
with carrier \(P\) is obtained by interpreting its input category in the universal operator
category:
\[
        \ulcorner\Phi\urcorner:A\to\mathsf U_B(P),
        \qquad
        \Phi=\ulcorner\Phi\urcorner^{\ast}\mathbb U_{B,P}
\]
as a fixed-carrier Mealy structure.  Equivalently, \(\mathbb U_{B,P}\) is terminal in the category
of fixed-carrier evaluators.

The one-step operational full-abstraction theorem says that equality of universal instructions is
exactly equality of evaluated transition-output effect on \(P\):
\[
        \varphi^{\operatorname{op}}=\psi^{\operatorname{op}}
        \quad\Longleftrightarrow\quad
        \operatorname{ev}(\varphi,-)=\operatorname{ev}(\psi,-)
\]
as represented sections over the relevant source-state fibre.

The extensional-collapse theorem says that any instruction syntax \(A\) maps into
\(\mathsf U_B(P)\) through its carrier-level operational quotient:
\[
        A
        \twoheadrightarrow
        \operatorname{CarOp}(\Phi)^{\operatorname{op}}
        \hookrightarrow
        \mathsf U_B(P).
\]

The intensional universal program category is
\[
        \mathsf{Prog}_B(P)
        =
        \mathrm{FreeCat}\bigl(\mathsf U_B(P)_{\mathrm{gr}}\bigr).
\]
Its arrows are formal finite programs in universal primitive instructions.  The canonical decoder
\[
        \operatorname{decode}_{B,P}:
        \mathsf{Prog}_B(P)\to\mathsf U_B(P)
\]
compiles such a formal program to its composite transition-output operator.  For a word
\[
        [\varphi_1^{\operatorname{op}},\ldots,\varphi_n^{\operatorname{op}}],
\]
the compiled operator is
\[
        (\varphi_1\diamond\cdots\diamond\varphi_n)^{\operatorname{op}},
\]
which executes \(\varphi_n\) first on the carrier, then \(\varphi_{n-1}\), and so on.

The corresponding intensional universal Mealy program machine is
\[
        \mathbb P_{B,P}
        =
        \operatorname{decode}_{B,P}^{\ast}\mathbb U_{B,P}.
\]

Thus the central architecture is
\[
\begin{tikzcd}[column sep=large,row sep=large]
\mathsf{Prog}_B(P)
\ar[r,"\operatorname{decode}_{B,P}"]
&
\mathsf U_B(P)
\ar[r,"\mathbb U_{B,P}"]
&
B .
\end{tikzcd}
\]
The first object is intensional.  The second is extensional.  The first map is compilation.  The
second Mealy morphism is evaluation.

The finite Mealy--Cayley theorem identifies the one-object finite-state operator category as
\[
        \mathrm{KlEnd}_N(Q)
        \cong
        N^Q\rtimes\operatorname{End}(Q),
\]
showing that the abstract construction recovers the full transformation-output monoid.

Finally, the Lawvere diagonal obstruction shows that full total extensional universality cannot
generally be self-coded by the same state object.  Classical universal computation therefore
requires additional intensional and partial structure: codes, partial evaluation, composition
tracking, parameterization, and iteration.

The resulting conclusion is:
\[
\boxed{
\begin{array}{c}
\text{Mealy morphisms may be intensional;}\\
\text{their universal operator semantics is extensional;}\\
\text{the canonical intensional layer is the free program category over that operator semantics.}
\end{array}}
\]